% =====================================================================
%  2bat-ud1-10.tex  —  SESSIÓ 10: Consolidació del producte: casos i errors
%  (sense preàmbul: s'inclou des de main.tex)
%  Criteri: C3 (producte de matrices). Consolidació abans del full de càlcul.
% =====================================================================
\horatitol{Sessió 10 — Consolidació del producte: casos i errors típics}%
{Assentar el producte de matrices amb una caça de l'error, casos difícils i la no-commutativitat.}

\subsection{Idea clau}

Abans de portar el producte de matrices al full de càlcul, val la pena
\textbf{assentar-lo bé} i conèixer els errors que es repeteixen. La regla d'or, com
sempre, és comprovar \emph{primer} la compatibilitat i preveure la dimensió del
resultat; després, fer el mecanisme fila-per-columna amb cura.

\begin{idea}
\textbf{Rutina i errors típics del producte}
\begin{itemize}[leftmargin=1.4em, itemsep=2pt]
  \item \textbf{Comprova la compatibilitat:} les \emph{columnes} de la primera han
        de coincidir amb les \emph{files} de la segona. Si no, el producte
        \textbf{no es pot fer}.
  \item \textbf{Preveu la dimensió:} una matriu $m\times n$ per una $n\times p$
        dona una matriu $m\times p$ (files de la primera, columnes de la segona).
  \item \textbf{Error 1:} multiplicar \emph{element a element} (com si fos una
        suma). El producte de matrices \emph{no} es fa així: és fila per columna.
  \item \textbf{Error 2:} intentar multiplicar matrices \emph{incompatibles}.
  \item \textbf{Error 3:} canviar l'\emph{ordre} dels factors pensant que dona el
        mateix. En general, $A\cdot B\neq B\cdot A$.
\end{itemize}
\end{idea}

\subsection{Exemples}

\begin{exemple}[caça l'error: element a element]
Algú ha escrit:
\[
\begin{pmatrix} 2 & 3 \\ 1 & 4 \end{pmatrix}\cdot\begin{pmatrix} 5 & 1 \\ 2 & 0 \end{pmatrix}
\overset{?}{=}\begin{pmatrix} 10 & 3 \\ 2 & 0 \end{pmatrix}.
\]
L'error és que ha multiplicat \emph{casella a casella} ($2\per 5$, $3\per 1$\dots),
com si fos una suma. El producte és fila per columna:
\[
\begin{pmatrix} 2 & 3 \\ 1 & 4 \end{pmatrix}\cdot\begin{pmatrix} 5 & 1 \\ 2 & 0 \end{pmatrix}
=\begin{pmatrix} 2\per 5+3\per 2 & 2\per 1+3\per 0 \\ 1\per 5+4\per 2 & 1\per 1+4\per 0 \end{pmatrix}
=\begin{pmatrix} 16 & 2 \\ 13 & 1 \end{pmatrix}.
\]
\end{exemple}

\begin{exemple}[caça l'error: matrices incompatibles]
Algú vol multiplicar una matriu $2\times 3$ per una $2\times 2$:
\[
\begin{pmatrix} 1 & 2 & 0 \\ 3 & 1 & 4 \end{pmatrix}\cdot\begin{pmatrix} 5 & 1 \\ 2 & 0 \end{pmatrix}.
\]
No es pot fer: les columnes de la primera ($3$) no coincideixen amb les files de
la segona ($2$). El producte \emph{no és compatible}.
\end{exemple}

\subsection{Exercicis resolts}

\begin{exresolt}[preveure la dimensió i calcular]
Calcula $A\cdot B$ amb
$A=\begin{pmatrix} 2 & 1 & 3 \end{pmatrix}$ (dimensió $1\times 3$) i
$B=\begin{pmatrix} 1 & 0 \\ 2 & 1 \\ 0 & 4 \end{pmatrix}$ (dimensió $3\times 2$).
\solucio
Primer comprovem: les columnes de $A$ ($3$) coincideixen amb les files de $B$
($3$): és compatible. La dimensió del resultat serà $1\times 2$.
\[
A\cdot B=\begin{pmatrix} 2\per 1+1\per 2+3\per 0 & 2\per 0+1\per 1+3\per 4 \end{pmatrix}
        =\begin{pmatrix} 4 & 13 \end{pmatrix}.
\]
\resposta{$A\cdot B=\begin{pmatrix} 4 & 13 \end{pmatrix}$ (dimensió $1\times 2$).}
\end{exresolt}

\begin{exresolt}[el repte: l'ordre importa]
Amb $A=\begin{pmatrix} 1 & 2 \\ 3 & 0 \end{pmatrix}$ i
$B=\begin{pmatrix} 0 & 1 \\ 2 & 1 \end{pmatrix}$, calcula $A\cdot B$ i $B\cdot A$
i comprova si surten iguals.
\solucio
\[
A\cdot B=\begin{pmatrix} 1\per 0+2\per 2 & 1\per 1+2\per 1 \\ 3\per 0+0\per 2 & 3\per 1+0\per 1 \end{pmatrix}
=\begin{pmatrix} 4 & 3 \\ 0 & 3 \end{pmatrix}.
\]
\[
B\cdot A=\begin{pmatrix} 0\per 1+1\per 3 & 0\per 2+1\per 0 \\ 2\per 1+1\per 3 & 2\per 2+1\per 0 \end{pmatrix}
=\begin{pmatrix} 3 & 0 \\ 5 & 4 \end{pmatrix}.
\]
No surten iguals: $A\cdot B\neq B\cdot A$. L'ordre dels factors importa.
\resposta{$A\cdot B=\begin{pmatrix} 4 & 3 \\ 0 & 3 \end{pmatrix}\neq B\cdot A=\begin{pmatrix} 3 & 0 \\ 5 & 4 \end{pmatrix}$: el producte no és commutatiu.}
\end{exresolt}

\subsection{Exercicis per fer}

\textit{Itinerari bàsic: exercicis 1 a 4. Ampliació: exercicis 5 a 7.}

\begin{exercicis}
\begin{exlist}
  \item Calcula $\begin{pmatrix} 3 & 1 \\ 2 & 4 \end{pmatrix}\cdot\begin{pmatrix} 1 & 2 \\ 0 & 3 \end{pmatrix}$.
  \item Calcula $\begin{pmatrix} 2 & 0 \\ 1 & 5 \end{pmatrix}\cdot\begin{pmatrix} 3 & 1 \\ 2 & 4 \end{pmatrix}$.
  \item \textbf{(Caça l'error)} Algú ha escrit
        \[ \begin{pmatrix} 1 & 2 \\ 3 & 1 \end{pmatrix}\cdot\begin{pmatrix} 2 & 0 \\ 1 & 4 \end{pmatrix}
           \overset{?}{=}\begin{pmatrix} 2 & 0 \\ 3 & 4 \end{pmatrix}. \]
        Explica quin error ha comès i escriu el producte correcte.
  \item Digues, \emph{sense calcular-lo}, quina dimensió tindrà el producte (o si
        no es pot fer):
        \begin{enumerate}[label=\alph*), leftmargin=1.6em, topsep=2pt]
          \item una $2\times 3$ per una $3\times 4$
          \item una $3\times 2$ per una $3\times 2$
          \item una $1\times 3$ per una $3\times 1$
        \end{enumerate}
  \item Calcula el producte rectangular
        $\begin{pmatrix} 1 & 2 & 1 \\ 0 & 1 & 3 \end{pmatrix}\cdot\begin{pmatrix} 2 & 1 \\ 1 & 0 \\ 0 & 2 \end{pmatrix}$
        i indica la dimensió del resultat.
  \item Amb $A=\begin{pmatrix} 2 & 1 \\ 0 & 3 \end{pmatrix}$ i
        $B=\begin{pmatrix} 1 & 4 \\ 2 & 0 \end{pmatrix}$, calcula $A\cdot B$ i
        $B\cdot A$ i comprova que són diferents.
  \item \textbf{(Raonament)} Tenim $A$ de dimensió $2\times 3$ i $B$ de dimensió
        $3\times 2$. Es poden fer tots dos productes? Quina dimensió té cadascun i
        què en dedueixes?
\end{exlist}
\end{exercicis}

% ---------------------------------------------------------------------
\fullsolucions{Sessió 10}
\begin{solucions}
\begin{sollist}
  \item $\begin{pmatrix} 3 & 9 \\ 2 & 16 \end{pmatrix}$.
  \item $\begin{pmatrix} 6 & 2 \\ 13 & 21 \end{pmatrix}$.
  \item L'error és haver multiplicat \emph{casella a casella} (com una suma). El
        producte correcte (fila per columna) és
        $\begin{pmatrix} 4 & 8 \\ 7 & 4 \end{pmatrix}$.
  \item (a) $2\times 4$ \quad (b) no es pot fer (columnes de la primera $=2$, files
        de la segona $=3$) \quad (c) $1\times 1$.
  \item $\begin{pmatrix} 4 & 3 \\ 1 & 6 \end{pmatrix}$ (dimensió $2\times 2$).
  \item $A\cdot B=\begin{pmatrix} 4 & 8 \\ 6 & 0 \end{pmatrix}$,
        \;$B\cdot A=\begin{pmatrix} 2 & 13 \\ 4 & 2 \end{pmatrix}$: són diferents.
  \item Sí, totes dues es poden fer: $A\cdot B$ és $2\times 2$ i $B\cdot A$ és
        $3\times 3$. Com que ni tan sols tenen la mateixa dimensió, queda clar que
        l'ordre dels factors importa i que, en general, $A\cdot B\neq B\cdot A$.
\end{sollist}
\end{solucions}
